\section{The datum and the finite colored evolutions}
\label{sec:construction}

The datum combines narrow jets with a bounded background close to a
Maxwellian.  The background supplies collisions at later times, whereas
an initial spatial hole keeps it apart from the jets during the short
interval on which we select the two families of solutions.  We first
construct these densities and their finite approximations.  The estimates
in this section and the next two sections hold for every fixed choice
\begin{equation}
 K\ge1,\qquad 0<\eta\le\log2,\qquad
 0<\alpha\le\tfrac14,\qquad 0<\gamma\le1.
 \label{con:parameters}
\end{equation}
The growth argument will fix these four parameters.  All remaining
small parameters are chosen after them.  The annular jets and the
colored evolution adapt the construction of
\cite[Sections 4--5]{OpenAIPeriodic}; the spatial orientation, the
background, and the later transfer of colors will allow uniform bounds
for arbitrarily large physical times.

\subsection{An irrational axis and annular jets}

Use the orthonormal frame
\[
 \frac{(-\sqrt2,1,0)}{\sqrt3},\qquad (0,0,1),\qquad
 \frac{(1,\sqrt2,0)}{\sqrt3}.
\]
In these coordinates write \(x=(y,z)\), \(v=(\bar v,v_3)\), and let
\(\Lambda\) be the image of \(\mathbb Z^3\) under the change of
coordinates.  The spatial torus is now \(\mathbb R^3/\Lambda\), with
covolume one.  Rotation preserves transport and the hard-sphere
collision operator, so this change does not alter the equation.

\begin{lemma}[Separation from the axis]\label{con:lattice}
There is a numerical constant \(c>0\) such that
\begin{equation}
 |\bar n|\ge\frac{c}{1+|n|}\qquad
             (n\in\Lambda\setminus\{0\}).
 \label{con:diophantine}
\end{equation}
\end{lemma}
\begin{proof}
An original integer vector \((k,m,\ell)\) has transverse coordinates
\(((m-\sqrt2 k)/\sqrt3,\ell)\).  If \(\ell\ne0\), its transverse
length is at least one.  Otherwise \((k,m)\ne(0,0)\), and
\[
 |m-\sqrt2 k|
 =\frac{|m^2-2k^2|}{|m+\sqrt2 k|}
 \ge\frac{1}{C(1+|k|+|m|)}.
\]
The numerator is a nonzero integer.  Since rotation preserves length,
this proves the claim.
\end{proof}

Fix a nonnegative \(J\in C_c^\infty(\mathbb R^3)\), supported in the
unit ball and of integral one.  With \(0<s_0\le1/8\), define
\begin{equation}
 s_j=s_0e^{-j\eta},\quad p_j=s_j^2,\quad
 c_\eta=\frac{e^\eta-1}{16},\quad
 d_s=d_0s^{10},\quad S(r)=\operatorname{sech}^2r,
 \qquad j\ge1,
 \label{con:scales}
\end{equation}
where \(0<d_0\le1\) is fixed sufficiently small.  It is enough to
make all Gaussian widths below, including their fixed widenings, fit
within the sheet estimates of Section~\ref{ker:section}.  Further
conditions involving positive powers of \(s_0\) will be obtained by
reducing \(s_0\).  Set
\[
 a_s^\sigma(z)=\sigma+\alpha\tanh(\gamma z/s),\qquad
 a=\tfrac18,\qquad b_{\mathrm{sp}}=\tfrac1{16},
 \qquad \sigma\in\{-1,+1\}.
\]
For each scale let \(E_j^\sigma\subset\mathbb R^2\) be the annulus
\[
 s_j+c_\eta p_j<|y|<e^\eta s_j-c_\eta p_j
\]
intersected with the stripes
\[
 y_1/p_j\pmod1\in
 \begin{cases}
 (0,1/4),&\sigma=+1,\\
 (1/2,3/4),&\sigma=-1.
 \end{cases}
\]
Both sets are nonempty for small \(s_0\): the radial interval is
nonempty, and a circle at an interior radius crosses both stripe
patterns because their period is \(s_j^2\).  Each pattern has spatial
fraction \(1/4\); no exact fraction is asserted after intersection
with an annulus.

A jet label is \(i=(j,\sigma,n)\), with \(n=(\bar n,n_3)\in\Lambda\).
For \(s=s_j\) and \((y_c,z_c)=x-n\), define
\begin{equation}
 c_i^0(x,v)=
 \ind_{E_j^\sigma}(y_c)\ind_{(-a,a)}(z_c)
 \frac K{s}S(\gamma z_c/s)d_s^{-3}
 J\!\left(\frac{v-a_s^\sigma(z_c)e_3}{d_s}\right).
 \label{con:jet-datum}
\end{equation}
All sums over labels include the lattice copies.  Translating \(x\)
by an element of \(\Lambda\) permutes these labels.

\subsection{The background and the initial entropy}

For any positive Maxwellian \(\rho\), write
\[
 \mathcal H_\rho(F)=\int
       \bigl(F\log(F/\rho)-F+\rho\bigr)\dd x\dd v.
\]
This agrees with the entropy already defined when \(\rho=M\).
Fix a sufficiently large numerical \(\beta_0\), and write
\(\mu(v)=e^{-\beta_0|v|^2}\).  Let \(0<\delta<1/32\), and let
\(D_{\mathrm{hole}}\) be the open \(2\delta\)-neighborhood in the
spatial torus of the axial segment of half-length \(1/3\) centered at
zero.  For a velocity cutoff \(R_\mu<\infty\), put
\begin{equation}
 u^0(x,v)=\ind_{D_{\mathrm{hole}}^c}(x)
             \mu(v)\ind_{\{|v|\le R_\mu\}},\qquad
 F_0=u^0+\sum_i c_i^0.
 \label{con:datum-equation}
\end{equation}
The hole has volume tending to zero as \(\delta\downarrow0\), by the
usual volume bound for a cylinder and its two endpoint neighborhoods.
For \(s_0\) sufficiently small compared with \(\delta\), every jet
lies in this hole.  The different jet supports are disjoint: annuli
and opposite stripe patterns are disjoint within a single copy, and
the unshifted supports lie in a ball of radius less than \(1/2\), so
distinct lattice copies cannot overlap.

Figure~\ref{fig:jet-geometry} summarizes the initial geometry.
\begin{figure}[tbp]
\centering
\resizebox{0.98\linewidth}{!}{%
\begin{tikzpicture}[x=1cm,y=1cm,>=Latex,font=\small,
  jetplus/.style={fill=blue!55!black,draw=none},
  jetminus/.style={fill=orange!75!black,draw=none},
  guide/.style={draw=black!40,line width=.35pt}]
  \path[use as bounding box] (0,0) rectangle (13.7,9.45);

  \node[anchor=west,font=\small\bfseries] at (0.15,9.15)
    {(a) Transverse annuli and stripes};
  \begin{scope}[shift={(3.1,6.65)}]
    \foreach \inner/\outer/\period in
      {1.025/1.475/.36,.678/0.989/.16,.449/0.661/.0711}{
      \begin{scope}[even odd rule]
        \clip (0,0) circle[radius=\outer]
          (0,0) circle[radius=\inner];
        \foreach \k in {-23,...,23}{
          \path[jetplus] ({\k*\period},-1.6)
            rectangle ({(\k+.25)*\period},1.6);
          \path[jetminus] ({(\k+.5)*\period},-1.6)
            rectangle ({(\k+.75)*\period},1.6);
        }
      \end{scope}
    }
    \foreach \radius in {1.5,1,.6667,.4444}
      \draw[guide] (0,0) circle[radius=\radius];
    \draw[->,black!65] (-1.85,0) -- (1.98,0) node[right] {$y_1$};
    \draw[->,black!65] (0,-1.8) -- (0,1.82) node[above] {$y_2$};
    \fill (0,0) circle[radius=.025];
    \node[anchor=east,fill=white,inner sep=1pt] at (-.06,-.15) {$0$};
    \draw[->,thin] (1.58,1.36) -- (1.13,.985);
    \node[anchor=west] at (1.53,1.48) {$e^\eta s_j$};
    \draw[->,thin] (1.78,.78) -- (.84,.54);
    \node[anchor=west] at (1.79,.83) {$s_j$};
  \end{scope}
  \node[align=center] at (3.1,4.63)
    {Geometric annuli; stripe period $p_j=s_j^2$.};

  \node[anchor=west,font=\small\bfseries] at (.15,3.76)
    {(b) A stripe period, magnified};
  \node[anchor=west] at (.25,3.28) {$y_1/p_j\pmod 1$};
  \path[jetplus] (.65,2.35) rectangle (1.9,2.92);
  \path[jetminus] (3.15,2.35) rectangle (4.4,2.92);
  \draw[black!60] (.65,2.35) rectangle (5.65,2.92);
  \foreach \xx in {1.9,3.15,4.4}
    \draw[black!45] (\xx,2.35)--(\xx,2.92);
  \node[text=white] at (1.275,2.635) {$+$};
  \node[text=white] at (3.775,2.635) {$-$};
  \node at (2.525,2.635) {empty};
  \node at (5.025,2.635) {empty};
  \foreach \xx/\lab in {.65/0,1.9/{\tfrac14},3.15/{\tfrac12},4.4/{\tfrac34},5.65/1}
    \node[below] at (\xx,2.30) {$\lab$};
  \node[align=center] at (3.15,1.12)
    {Each sign occupies one quarter\\
     of every stripe period.};

  \node[anchor=west,font=\small\bfseries] at (7.05,9.15)
    {(c) Jets inside the bath hole};
  \path[fill=black!9] (7.1,5.05) rectangle (13.55,8.15);
  \path[fill=white,draw=black!55]
    (8.8,5.90) -- (11.8,5.90)
    arc[start angle=-90,end angle=90,radius=.70]
    -- (8.8,7.30)
    arc[start angle=90,end angle=270,radius=.70] -- cycle;
  \draw[->,black!50,dashed] (7.4,6.60) -- (13.25,6.60)
    node[above] {$z$};
  \draw[black!45,line width=1pt] (8.8,6.60)--(11.8,6.60);
  \foreach \xx/\lab in {8.8/{-\tfrac13},10.3/0,11.8/{\tfrac13}}
    {\draw[black!60] (\xx,6.51)--(\xx,6.69);
     \node[below,fill=white,inner sep=1pt] at (\xx,6.49) {$\lab$};}
  \path[jetplus] (9.7375,6.85) rectangle (10.8625,6.94);
  \path[jetminus] (9.7375,6.98) rectangle (10.8625,7.07);
  \draw[<->,thin] (9.7375,7.20)--(10.8625,7.20)
    node[midway,above,inner sep=1pt] {$|z|<a$};
  \draw[->,thin] (12.65,7.58)--(12.31,7.02);
  \node[anchor=east] at (13.40,7.81) {$D_{\rm hole}$};
  \node[anchor=west,align=left] at (7.3,5.45)
    {Initial bath on $D_{\rm hole}^{\,c}$};
  \node[align=center] at (10.3,4.63)
    {Schematic axial view; $a=1/8$.};

  \node[anchor=west,font=\small\bfseries] at (7.05,3.76)
    {(d) Axial velocity centers};
  \draw[->,black!55] (7.7,1.87)--(13.30,1.87) node[right] {$z$};
  \draw[->,black!55] (10.3,.30)--(10.3,3.23) node[above] {$v_3$};
  \draw[dotted,black!45] (7.9,2.55)--(12.7,2.55);
  \draw[dotted,black!45] (7.9,1.19)--(12.7,1.19);
  \node[anchor=east,fill=white,inner sep=1pt] at (10.22,2.55) {$1$};
  \node[anchor=east,fill=white,inner sep=1pt] at (10.22,1.19) {$-1$};
  \draw[blue!55!black,line width=1.1pt,domain=-2.2:2.2,samples=70]
    plot ({10.3+\x},{2.55+.23*tanh(1.35*\x)});
  \draw[orange!75!black,line width=1.1pt,domain=-2.2:2.2,samples=70]
    plot ({10.3+\x},{1.19+.23*tanh(1.35*\x)});
  \node[anchor=west,text=blue!55!black,fill=white,inner sep=1pt]
    at (12.55,2.78) {$a_s^+$};
  \node[anchor=west,text=orange!75!black,fill=white,inner sep=1pt]
    at (12.55,1.40) {$a_s^-$};
  \node[align=center] at (10.3,.04)
    {$a_s^\sigma(z)=\sigma+\alpha\tanh(\gamma z/s)$};
\end{tikzpicture}%
}
\caption{The initial jets and bath near one lattice point.
The annular picture is schematic: stripe widths and radial gaps are enlarged
for visibility, and only three annuli are drawn.  Panel (b) records the exact
stripe fractions, which later give a coefficient $1/4$ for each sign after
averaging.  Panel (c) places the jets schematically inside the bath hole,
the tube of radius $2\delta$ about the axial segment $[-1/3,1/3]$;
the vertical placement of the two signs is only illustrative.
Panel (d) shows the centers of the compact velocity profiles of width $d_s$.
All spatial panels are local views in the rotated lift; they do not depict
a closed axis on the torus.}
\label{fig:jet-geometry}
\end{figure}


Choose also a nonzero nonnegative
\(\chi\in C_c^\infty(\mathbb R^3_x\times\mathbb R^3_v)\), and a
constant \(L_0\) with \(L_0\|\chi\|_H>1\).  For the approximation
retaining \(j\le N\), the additional initial density is either zero
or
\begin{equation}
 h^0_{N,T,\lambda}(x,v)
 =\lambda m_Ts_N\sum_{n\in\Lambda}
        \chi\bigl((x-n)/s_N,v\bigr),\qquad
 m_T=T^p,\qquad 0\le\lambda\le L_0.
 \label{con:seed}
\end{equation}
Here \(0<T\le t_1\le1/32\); Section~\ref{sec:early} fixes
\(t_1>0\) and then \(p>0\).  We reduce \(s_0\) so that the seed's
lattice copies have disjoint spatial supports.  Its weighted norm is
then \(\lambda m_Ts_N\|\chi\|_H\), with the unperiodized norm on the
right.

\begin{lemma}[Datum bounds and approximation]\label{con:datum}
The density \(F_0\) is nonnegative, has bounded velocity support, and
satisfies
\[
 \int F_0(1+|v|^2+|\log F_0|)\dd x\dd v<\infty.
\]
Its truncations with \(j\le N\), with or without a seed
\eqref{con:seed}, converge to \(F_0\) in weighted \(L^1\) for every
fixed polynomial velocity weight and in relative entropy with respect
to either \(M\) or \(\mu\).  The convergence is uniform for
\(0<T\le t_1\) and \(0\le\lambda\le L_0\).
Moreover their initial \(\mu\)-relative entropy can be made arbitrarily
small, uniformly in \(N,T,\lambda\), by first decreasing \(\delta\)
and increasing \(R_\mu\), then decreasing \(s_0\).
\end{lemma}
\begin{proof}
Integration over one period of the lattice sum equals integration of
each unshifted prototype over \(\mathbb R^3_x\).  The transverse area
is \(O(s^2)\), while
\[
 \int_{-a}^a S(\gamma z/s)\dd z\le 2s/\gamma.
\]
The velocity factor has integral one and support in \(\{|v|<2\}\).
Thus every fixed velocity moment of a sign prototype is \(O(s^2)\).
On its positive set the logarithm of the density splits into
\[
 \log((K/s)d_s^{-3}),\qquad \log S(\gamma z/s),\qquad \log J(w).
\]
The first term is \(O(1+|\log s|)\), and
\(\int S|\log S|<\infty\), \(\int J|\log J|<\infty\).
Its absolute entropy is therefore
\(O(s^2(1+|\log s|))\).  The corresponding series over the geometric
scales converges.  Disjointness proves the asserted integrability and
convergence for the cold sum.  The background is bounded by \(\mu\),
and its velocities satisfy \(|v|\le R_\mu\), giving the datum claims.

For a seed \(b=h^0_{N,T,\lambda}\), scaling \(x\) gives
\[
 \int b\langle v\rangle^k\dd x\dd v\le C_k s_N^4,\qquad
 \int b|\log b|\dd x\dd v
       \le C s_N^4(1+|\log s_N|).
\]
These constants are uniform because \(\lambda m_T\) stays in a fixed
bounded interval and \(r|\log r|\) is bounded on such an interval.
The preexisting truncated density \(f=u^0+\sum_{j\le N}c_i^0\)
satisfies \(f\le C s_N^{-31}\).  For \(f,b\ge0\), integration of
\(1+\log(f+r)\) over \(0<r<b\) gives
\[
 |(f+b)\log(f+b)-f\log f|
 \le b\{2+\log^+(f+b)\}+b|\log b|.
\]
Indeed, the positive logarithm is bounded by its terminal value, while
\(\int_0^b\log^-(f+r)\dd r\le b(1+|\log b|)\).
The entropy change caused by the seed is consequently
\(O(s_N^4(1+|\log s_N|))\).  Its reference terms are controlled by
its mass and energy.  This proves all convergence assertions.

Removing the hole and the velocity tail from \(\mu\) costs precisely
the removed mass in \(\mathcal H_\mu\).  The inserted jets have
bounded velocities, so their contribution is bounded by the preceding
absolute entropy, mass, and energy sums, which tend to zero with
\(s_0\).  The same is true uniformly for the seeds.  The stated order
of choices therefore makes the initial relative entropy arbitrarily
small.
\end{proof}

\subsection{Colored equations and transfers at large times}

Use the collision cutoff obtained by multiplying the hard-sphere
kernel by
\[
 \vartheta_b(|v-v_*|)=\min\{1,b/|v-v_*|\},\qquad b\ge1.
\]
Its value at zero is immaterial.  Write \(P^b\) for the symmetric
cutoff gain and
\(\nu_f^b(v)=2\pi\int\min\{b,|v-p|\}f(p)\dd p\).
The cutoff is invariant under particle exchange and the elastic
pre/post-collisional substitution.  It can equivalently be inserted
in the incoming relative speed in the Carleman and sphere formulas.
Fix \(b=b_N\to\infty\), and put
\(\mathcal I_N=\{(j,\sigma,n):1\le j\le N,\ \sigma=\pm1,
 n\in\Lambda\}\).

Until \(t_1\), write \(c=\sum_{i\in\mathcal I_N}c_i\),
\(F=c+u+h\), and solve
\begin{equation}
 \begin{aligned}
 \D c_i+\nu_F^b c_i&=P^b(c_i,c_i),&
 \D u+\nu_F^b u&=P^b(u,u),\\
 \D h+\nu_F^b h&=P^b(h,h)+2P^b(c+u,h)+\mathcal S,&
 \mathcal S&=2P^b(c,u)+\sum_{i\ne i'}P^b(c_i,c_{i'}).
 \end{aligned}
 \label{con:colors}
\end{equation}
The last sum is ordered.  Thus \(\D F=Q^b(F,F)\).

At \(t_1\), replace \(u+h\) by a single component \(w\).  A jet at
scale \(s\) remains separate until its retirement time
\begin{equation}
 L_s=s^{-1/2}>1.
 \label{con:deadline}
\end{equation}
At \(L_s\), add all colors at that scale to \(w\) and remove their
labels.  For the active labels, retain their equations in
\eqref{con:colors}; between retirements set \(F=c+w\) and solve
\begin{equation}
 \D w+\nu_F^b w
   =P^b(w,w)+2P^b(c,w)+\sum_{i\ne i'}P^b(c_i,c_{i'}).
 \label{con:late-colors}
\end{equation}
The total density has no jump at a retirement.  It satisfies the same
cutoff Boltzmann equation across the entire evolution.

\begin{proposition}[Finite evolution and its identities]
\label{con:finite-evolution}
For fixed \(N,b,T,\lambda\), the colored equations have a unique
nonnegative local mild solution with bounded total \(H\)-norm.  It
continues while that norm remains bounded, with the prescribed
transfers at retirement times, and depends continuously in this norm
on its initial data on every common bounded interval of existence.
The norm of each component is continuous between transfers.

On every such interval, total mass, momentum, and energy are conserved.
The exact local collision-invariant balances hold, with all polynomial
velocity fluxes integrable.  For \(\rho=M\) or \(\rho=\mu\), one has
at every time of existence
\begin{equation}
 \mathcal H_\rho(F(t))+
       \int_0^t\mathcal D^b(F(s))\dd s
       \le\mathcal H_\rho(F(0)),
 \label{con:finite-entropy}
\end{equation}
where \(\mathcal D^b\) is the entropy dissipation with the same cutoff.
\end{proposition}
\begin{proof}
The gain estimate \(P(H,H)\le CH\) and
\[
 \|\nu_f^b\|_\infty\le
  2\pi b\|f\|_H\int H(p)\dd p
\]
make the gain and loss locally Lipschitz in the weighted norm.
For the colored system use the norm of the pointwise sum of the
absolute values of the components, and the sum of absolute differences
for two solutions.  All quadratic estimates close in this norm.
The initially disjoint lattice copies have finite total norm, and
periodic covariance is preserved by uniqueness.  Picard iteration in
free coordinates gives the local solution; the loss integrating factor
and positive gain iteration give nonnegativity.  Repetition on short
subintervals gives continuation and the difference estimate.  These
arguments use the integral equation along flights and require no norm
continuity of spatial translation for measurable data.  On an interval
with bounded norm, the increments along flights are Lipschitz in that
norm.  Since free transport is isometric, each component's norm is
continuous.  Transfers only regroup finitely many scales and preserve
the total density.

The weight \(H\) has every polynomial moment and belongs to
\(L^{1+\epsilon}(\mathbb R^3)\) for some \(\epsilon>0\).  Hence the
cutoff equation has integrable gain and loss with all polynomial
weights on every bounded interval under consideration.  Integrating
its path formula and using particle interchange and the elastic
substitution yields conservation of the collision invariants.  The
same calculation with a smooth spatial test gives each local balance,
including its initial trace.

Here are details of the entropy calculation at a fixed approximation.
Put \(g(v)=e^{-2|v|^2}\) and \(F_\epsilon=F+\epsilon g\).  Differentiate
\(F_\epsilon\log(F_\epsilon/\rho)-F_\epsilon+\rho\) along flights;
its derivative is \(Q^b(F,F)\log(F_\epsilon/\rho)\).
Replacing \(Q^b(F,F)\) by \(Q^b(F_\epsilon,F_\epsilon)\) changes its
integral by a quantity tending to zero.  Indeed, the operator difference
is \(O(\epsilon)H\) on the fixed interval, while
\[
 |\log(F_\epsilon/\rho)|
 \le C\bigl(1+|\log\epsilon|+|v|^2+
                           \log^+(1/|\bar v|)\bigr).
\]
The right-hand side is integrable against \(H\), and the resulting
error is \(O(\epsilon(1+|\log\epsilon|))\).
For the shifted density the usual symmetric calculation is now
justified by this integrability.  The logarithm of either Maxwellian
reference is a collision invariant.  Symmetrization therefore gives
minus the cutoff dissipation.  The regularized entropies converge by
dominated convergence using the same bounds, and Fatou's lemma applies
to the nonnegative dissipation.  This proves
\eqref{con:finite-entropy}.  The integrated path formula includes every
time section, so the conservation and entropy statements hold at every
time, also when colors are transferred.
\end{proof}

We record the consequence of the small initial entropy that will drive
the later estimates.  Define
\begin{equation}
 D_F(t,x)=\int\langle v\rangle|F(t,x,v)-\mu(v)|\dd v,\qquad
 D_-(t,x)=\int\langle v\rangle(\mu(v)-F(t,x,v))_+\dd v.
 \label{con:entropy-defect}
\end{equation}
For any prescribed sufficiently small \(\varepsilon_E>0\), the
parameters in Lemma~\ref{con:datum} can be chosen so that, throughout
every interval of existence of every finite approximation,
\begin{equation}
 \|D_F(t)\|_{L^1_x}\le\varepsilon_E,\qquad
 0\le D_-(t,x)\le\int\langle v\rangle\mu(v)\dd v.
 \label{con:entropy-distance}
\end{equation}
To see the first assertion, use
\((\sqrt z-1)^2\le z\log z-z+1\) and Cauchy--Schwarz to obtain
\[
 \|D_F(t)\|_{L^1_x}
 \le\mathcal H_\mu(F(t))^{1/2}
       \left(2\int\langle v\rangle^2(F(t)+\mu)\dd x\dd v\right)^{1/2}.
\]
The second factor has a uniform fixed bound by mass and energy
conservation and the initial estimates.  The first factor is uniformly
small by \eqref{con:finite-entropy}.  No pointwise smallness of \(D_F\)
is asserted here; recovering the necessary pointwise control after
transport is the purpose of the global continuation argument.
